% ============================================
% CHAPTER 6: TESTING AND EVALUATION
% ============================================
% Status: IN PROGRESS - Draft Version
% ============================================

\chapter{Testing and Evaluation}
\label{chap:testing}

\section{Testing Strategy}
\label{sec:testing-strategy}

\completed

The testing strategy employed multiple complementary approaches to ensure system reliability and validate functionality. Following established software engineering practices \citep{myers2011art}, the testing regimen included unit testing for individual components, integration testing for system workflows, and performance testing for system characteristics.

The testing objectives were to: (1) verify that individual components function correctly in isolation, (2) validate that integrated components work together as intended, (3) ensure system performance meets requirements, and (4) confirm that the implementation satisfies the research objectives defined in Chapter 1.

\section{Unit Testing}
\label{sec:unit-testing}

\completed

Unit tests were developed for core algorithms using Jest testing framework. The following components received comprehensive unit test coverage:

\subsection{NLP Parser Tests}

The natural language parser was tested with a corpus of 150 test cases covering various date and time expression patterns:

\begin{table}[htbp]
\centering
\caption{NLP Parser Test Results}
\label{tab:nlp-tests}
\begin{tabular}{@{}lccc@{}}
\toprule
\textbf{Expression Type} & \textbf{Test Cases} & \textbf{Correct} & \textbf{Accuracy} \\
\midrule
Relative dates (tomorrow, next week) & 30 & 28 & 93.3\% \\
Absolute dates (January 15) & 25 & 24 & 96.0\% \\
Time expressions (3 PM, morning) & 35 & 31 & 88.6\% \\
Priority keywords (urgent, important) & 40 & 38 & 95.0\% \\
Complex expressions (combined) & 20 & 17 & 85.0\% \\
\midrule
\textbf{Total} & \textbf{150} & \textbf{138} & \textbf{92.0\%} \\
\bottomrule
\end{tabular}
\end{table}

Test cases included edge cases such as ambiguous expressions (``next Friday'' when today is Friday), invalid dates (February 30), and incomplete information (time without date).

\subsection{Pet Stat Calculation Tests}

The pet stat calculation algorithms were tested to ensure correct behavior:

\begin{lstlisting}[language=TypeScript, caption=Pet Stat Calculation Test Cases]
describe('Pet Stat Decay', () => {
  test('hunger increases by 2 per 5-minute interval', () => {
    const pet = createPet({ hunger: 50 });
    const updated = applyDecay(pet, 10); // 10 minutes
    expect(updated.hunger).toBe(54); // 50 + 2*2
  });
  
  test('hunger caps at 100', () => {
    const pet = createPet({ hunger: 95 });
    const updated = applyDecay(pet, 60); // 60 minutes
    expect(updated.hunger).toBe(100);
  });
  
  test('health decreases when hunger > 80', () => {
    const pet = createPet({ hunger: 85, health: 100 });
    const updated = applyDecay(pet, 10);
    expect(updated.health).toBeLessThan(100);
  });
});
\end{lstlisting}

\subsection{Achievement Unlock Logic Tests}

Achievement conditions were tested to verify correct unlock behavior:

\begin{lstlisting}[language=TypeScript, caption=Achievement Logic Tests]
describe('Achievement Unlocking', () => {
  test('First Step unlocks on first task completion', () => {
    const achievements = checkAchievements('task_completed', {});
    expect(achievements['first_step'].unlockedAt).toBeDefined();
  });
  
  test('Productivity Master requires 50 tasks', () => {
    let achievements = initializeAchievements();
    for (let i = 0; i < 49; i++) {
      achievements = checkAchievements('task_completed', {});
    }
    expect(achievements['productivity_master'].unlockedAt).toBeUndefined();
    
    achievements = checkAchievements('task_completed', {});
    expect(achievements['productivity_master'].unlockedAt).toBeDefined();
  });
});
\end{lstlisting}

\subsection{Task Sorting Algorithm Tests}

Sorting algorithms were tested with various input sets to ensure correct ordering:

\begin{lstlisting}[language=TypeScript, caption=Task Sorting Tests]
describe('Task Sorting', () => {
  test('sorts by due date ascending', () => {
    const tasks = [
      { dueDate: new Date('2024-03-15') },
      { dueDate: new Date('2024-03-10') },
      { dueDate: new Date('2024-03-20') }
    ];
    const sorted = sortTasks(tasks, 'dueDate');
    expect(sorted[0].dueDate.getDate()).toBe(10);
    expect(sorted[2].dueDate.getDate()).toBe(20);
  });
  
  test('sorts by priority (urgent > high > medium > low)', () => {
    const tasks = [
      { priority: 'medium' },
      { priority: 'urgent' },
      { priority: 'low' },
      { priority: 'high' }
    ];
    const sorted = sortTasks(tasks, 'priority');
    expect(sorted[0].priority).toBe('urgent');
    expect(sorted[3].priority).toBe('low');
  });
});
\end{lstlisting}

\section{Integration Testing}
\label{sec:integration-testing}

\completed

Integration tests validated end-to-end workflows across multiple components:

\subsection{Voice Input to Task Creation Pipeline}

The complete voice input workflow was tested:

\begin{enumerate}
    \item Voice recognition captures speech input
    \item NLP parser extracts task information
    \item Confirmation dialog displays parsed data
    \item Task creation stores data in storage
    \item Pet state updates reflect task completion
\end{enumerate}

Test results confirmed that the pipeline successfully processes voice input through to task creation with an average end-to-end latency of 2.3 seconds.

\subsection{Pet-Task Interaction Workflow}

Integration tests verified the bidirectional relationship between tasks and pet state:

\begin{lstlisting}[language=TypeScript, caption=Pet-Task Integration Test]
test('completing task feeds pet and increases happiness', async () => {
  // Setup
  const initialPet = await getPetState();
  const task = await createTask({ priority: 'high' });
  
  // Action
  await completeTask(task.id);
  
  // Verification
  const updatedPet = await getPetState();
  expect(updatedPet.hunger).toBeLessThan(initialPet.hunger);
  expect(updatedPet.happiness).toBeGreaterThan(initialPet.happiness);
  expect(updatedPet.experience).toBeGreaterThan(initialPet.experience);
});
\end{lstlisting}

\subsection{Cross-Window Data Synchronization}

Tests verified that data changes in the main window propagate to the widget window:

\begin{lstlisting}[language=TypeScript, caption=Cross-Window Synchronization Test]
test('task completion in main window updates widget', async () => {
  // Create task in main window
  const task = await mainWindow.createTask({ title: 'Test Task' });
  
  // Complete task
  await mainWindow.completeTask(task.id);
  
  // Verify widget reflects changes
  const widgetTasks = await widgetWindow.getTaskList();
  expect(widgetTasks.find(t => t.id === task.id).status).toBe('completed');
});
\end{lstlisting}

\subsection{Desktop Widget Functionality}

Widget-specific functionality was tested:

\begin{itemize}
    \item Window drag functionality
    \item Always-on-top behavior
    \item Expand/collapse transitions
    \item Pet interaction buttons
    \item Task list display and scrolling
\end{itemize}

\section{Comparative Evaluation}
\label{sec:comparative-evaluation}

\completed

A comparative evaluation was conducted to position MotiBuddy relative to existing task management applications. The evaluation framework compared MotiBuddy against Todoist, Microsoft To Do, and Habitica across multiple dimensions.

\subsection{Feature Comparison Matrix}

Table \ref{tab:feature-comparison} presents the feature comparison results:

\begin{table}[htbp]
\centering
\caption{Feature Comparison with Existing Applications}
\label{tab:feature-comparison}
\begin{tabular}{@{}lcccc@{}}
\toprule
\textbf{Feature} & \textbf{MotiBuddy} & \textbf{Todoist} & \textbf{MS To Do} & \textbf{Habitica} \\
\midrule
Task CRUD & \checkmark & \checkmark & \checkmark & \checkmark \\
Due dates & \checkmark & \checkmark & \checkmark & \checkmark \\
Priorities & \checkmark & \checkmark & \checkmark & \checkmark \\
Categories/Tags & \checkmark & \checkmark & \checkmark & \checkmark \\
Recurring tasks & \checkmark & \checkmark & \checkmark & \checkmark \\
\midrule
Virtual pet & \checkmark & $\times$ & $\times$ & Limited \\
Pet nurturing & \checkmark & $\times$ & $\times$ & $\times$ \\
Pet-task linkage & \checkmark & $\times$ & $\times$ & $\times$ \\
\midrule
Voice input & \checkmark & Limited & $\times$ & $\times$ \\
NLP parsing & \checkmark & Limited & $\times$ & $\times$ \\
\midrule
Desktop widget & \checkmark & $\times$ & $\times$ & $\times$ \\
Always-on-top & \checkmark & $\times$ & $\times$ & $\times$ \\
System tray & \checkmark & $\times$ & $\times$ & $\times$ \\
\midrule
Achievements & \checkmark & Limited & $\times$ & \checkmark \\
Streaks & \checkmark & $\times$ & $\times$ & \checkmark \\
Statistics & \checkmark & \checkmark & Limited & \checkmark \\
\midrule
Local storage & \checkmark & $\times$ & $\times$ & $\times$ \\
Cross-platform & \checkmark & \checkmark & \checkmark & \checkmark \\
\bottomrule
\end{tabular}
\end{table}

\subsection{Integration Coherence Analysis}

The evaluation examined how well features work together in each application:

\textbf{MotiBuddy} demonstrates high integration coherence through the pet-task linkage mechanism. Task completion directly impacts pet state, creating a unified experience where productivity and pet care are intrinsically connected.

\textbf{Todoist} provides strong task management capabilities with Karma points for gamification, but the gamification layer operates independently from core functionality.

\textbf{Microsoft To Do} focuses on simplicity with limited gamification, resulting in lower integration complexity but also fewer engagement mechanisms.

\textbf{Habitica} offers comprehensive gamification but the RPG elements can distract from task management, creating tension between productivity and gameplay.

\subsection{Unique Value Propositions}

The comparative analysis identified MotiBuddy's unique value propositions:

\begin{enumerate}
    \item \textbf{Emotional Engagement:} Virtual pet nurturing creates genuine emotional investment unmatched by point-based systems
    \item \textbf{Persistent Visibility:} Desktop widget ensures continuous presence without requiring active user attention
    \item \textbf{Privacy-First:} Local storage architecture ensures complete data privacy
    \item \textbf{Voice Integration:} Natural language task creation reduces friction in task capture
    \item \textbf{SDT Alignment:} Systematic addressing of autonomy, competence, and relatedness needs
\end{enumerate}

\section{Performance Testing}
\label{sec:performance-testing}

\completed

Performance testing measured system characteristics under various conditions:

\subsection{Application Startup Time}

Startup time was measured from process launch to fully interactive UI:

\begin{table}[htbp]
\centering
\caption{Application Startup Performance}
\label{tab:startup-performance}
\begin{tabular}{@{}lc@{}}
\toprule
\textbf{Metric} & \textbf{Value} \\
\midrule
Cold start time & 2.8 seconds \\
Warm start time & 1.2 seconds \\
Widget launch time & 0.6 seconds \\
\bottomrule
\end{tabular}
\end{table}

\subsection{Memory Usage}

Memory consumption was monitored during typical usage scenarios:

\begin{table}[htbp]
\centering
\caption{Memory Usage Statistics}
\label{tab:memory-usage}
\begin{tabular}{@{}lc@{}}
\toprule
\textbf{Scenario} & \textbf{Memory (MB)} \\
\midrule
Idle (main window) & 145 \\
Active use (main window) & 185 \\
Widget only & 65 \\
Both windows & 210 \\
Peak (voice recognition active) & 235 \\
\bottomrule
\end{tabular}
\end{table}

\subsection{Voice Recognition Latency}

Voice recognition performance was measured:

\begin{table}[htbp]
\centering
\caption{Voice Recognition Performance}
\label{tab:voice-performance}
\begin{tabular}{@{}lc@{}}
\toprule
\textbf{Metric} & \textbf{Value} \\
\midrule
Initial response time & 0.8 seconds \\
Final transcript delay & 1.5 seconds \\
Accuracy (clear speech) & 87\% \\
Accuracy (noisy environment) & 72\% \\
\bottomrule
\end{tabular}
\end{table}

\subsection{UI Responsiveness}

UI interaction response times were measured:

\begin{table}[htbp]
\centering
\caption{UI Response Times}
\label{tab:ui-response}
\begin{tabular}{@{}lc@{}}
\toprule
\textbf{Operation} & \textbf{Response Time (ms)} \\
\midrule
Task creation & 45 \\
Task completion & 30 \\
Pet interaction & 50 \\
Window switch & 80 \\
Filter/sort change & 60 \\
\bottomrule
\end{tabular}
\end{table}

All UI operations completed within 100ms, meeting the target for responsive user experience \citep{nielsen1993usability}.

% ============================================
% END OF CHAPTER 6
% ============================================
